\title{An integrated bioinformatics platform for investigating the human E3 ubiquitin ligase-substrate interaction network}
\begin{document}
\maketitle

\begin{abstract}
ARTICLE An integrated bioinformatics platform for investigating the human E3 ubiquitin ligase- substrate interaction network Yang Li1, Ping Xie1,2, Liang Lu1, Jian Wang1, Lihong Diao1,3, Zhongyang Liu1, Feifei Guo1, Yangzhige He1, Yuan Liu1, Qin Huang3, Han Liang 4, Dong Li1 \& Fuchu He1 The ubiquitination mediated by ubiquitin activating enzyme (E1), ubiquitin conjugating enzyme (E2), and ubiquitin ligase (E3) cascade is crucial to protein degradation, transcription regulation, and cell signaling in eukaryotic cells. The high specificity of ubiquitination is regulated by the interaction between E3 ubiquitin ligases and their target substrates. Unfortunately, the landscape of human E3-substrate network has not been systematically uncovered. Therefore, there is an urgent need to develop a high-throughput and efficient strategy to identify the E3-substrate interaction. To address this challenge, we develop a computational model based on multiple types of heterogeneous biological evidence to investigate the human E3-substrate interactions. Furthermore, we provide UbiBrowser as an integrated bioinformatics platform to predict and present the proteome-wide human E3-substrate interaction network (http://ubibrowser.ncpsb.org). DOI: 10.1038/s41467-017-00299-9 OPEN 1 State Key Laboratory of Proteomics, Beijing Proteome Research Center, Beijing Institute of Radiation Medicine, National Center for Protein Sciences (The PHOENIX Center, Beijing), Beijing 102206, China. 2 Department of Bioche
\end{abstract}

\section{Recovered Text}
ARTICLE
An integrated bioinformatics platform for
investigating the human E3 ubiquitin ligase-
substrate interaction network
Yang Li1, Ping Xie1,2, Liang Lu1, Jian Wang1, Lihong Diao1,3, Zhongyang Liu1, Feifei Guo1, Yangzhige He1, Yuan Liu1,
Qin Huang3, Han Liang
4, Dong Li1 \& Fuchu He1
The ubiquitination mediated by ubiquitin activating enzyme (E1), ubiquitin conjugating
enzyme (E2), and ubiquitin ligase (E3) cascade is crucial to protein degradation, transcription
regulation, and cell signaling in eukaryotic cells. The high specificity of ubiquitination is
regulated by the interaction between E3 ubiquitin ligases and their target substrates.
Unfortunately, the landscape of human E3-substrate network has not been systematically
uncovered. Therefore, there is an urgent need to develop a high-throughput and efficient
strategy to identify the E3-substrate interaction. To address this challenge, we develop a
computational model based on multiple types of heterogeneous biological evidence to
investigate the human E3-substrate interactions. Furthermore, we provide UbiBrowser as an
integrated bioinformatics platform to predict and present the proteome-wide human
E3-substrate interaction network (http://ubibrowser.ncpsb.org).
DOI: 10.1038/s41467-017-00299-9
OPEN
1 State Key Laboratory of Proteomics, Beijing Proteome Research Center, Beijing Institute of Radiation Medicine, National Center for Protein Sciences (The
PHOENIX Center, Beijing), Beijing 102206, China. 2 Department of Biochemistry and Molecular Biology, Capital Medical University, Beijing 100069, China.
3 School of Chemistry and Chemical Engineering, Guangxi University for Nationalities, Nanning 530006, China. 4 Department of Bioinformatics and
Computational Biology, The University of Texas MD Anderson Cancer Center, Houston, TX 77030, USA. Yang Li, Ping Xie and Liang Lu contributed equally
to this work. Correspondence and requests for materials should be addressed to D.L. (email: lidong.bprc@foxmail.com) or to F.H. (email: hefc@nic.bmi.ac.cn)
NATURE COMMUNICATIONS| 8:  347
| DOI: 10.1038/s41467-017-00299-9| www.nature.com/naturecommunications
1
U
biquitin, which is an abundant 76 amino acid polypeptide,
can covalently conjugate to certain proteins by an
isopeptide bond between its carboxyl and the amino
group of a lysine residue. This process is mediated by a cascade
consistent with ubiquitin activating enzyme (E1), ubiquitin con-
jugating enzyme (E2), and ubiquitin ligase (E3)1. Ubiquitination
regulates a wide spectrum of proteolytic and nonproteolytic cel-
lular processes in eukaryotes, including proteasome-mediated
protein degradation, inflammatory signaling, DNA damage
response, and enzymatic activity regulation2. Thus, ubiquitination
is closely related to the development of many diseases like Alz-
heimer’s disease3, 4, Parkinson’s disease5, and multiple cancers6–8.
During the ubiquitination procedure, the interaction between
ubiquitin ligase and substrate decides the substrate’s specificity
and
determines
the
substrate’s
fate.
Multiple
traditional
experimental strategies (e.g., GPS profiling9, protein micro-
arrays10, live phage display library11, and mass spectrometry12)
have been developed to identify the E3-substrate interaction
(ESI). However, because of E3s’ low substrate levels and their
intrinsically weak interactions with substrates, these methods are
laborious, time intensive, expensive, and low efficient. As a result,
although there are more than 30000 ubiquitin sites on over
5700 substrates in the ubiquitination site database mUbiSiDa
(version 1.0)13, only less than 900 human E3-substrate relation-
ships are collected in database14, which means that only a small
proportion (~15\%) of ubiquitinated proteins has the known
corresponding ubiquitin ligase information. Therefore, a robust
computational strategy is desirable to systematically identify the
potential E3-substrate interaction at proteome scale.
To
address
this
issue,
we
developed
a
naïve
Bayesian
classifier-based computational algorithm to combine multiple types
of heterogeneous biological evidence including homology E3-
substrate interaction, enriched domain and Gene Ontology (GO)
term pair, protein interaction network loop, and inferred E3
recognition consensus motif, to predict human E3-substrate inter-
actions. Then we implemented a proteome-wide E3-substrate
interactions scanning generating a predicted E3-substrate interac-
tion data set (PESID). Finally, to facilitate the usage of our algo-
rithm and PESID, we presented an online platform (UbiBrowser) to
investigate human ubiquitin ligase-substrate interaction network.
Results
Overview of our prediction protocol. UbiBrowser was designed
to be a naïve Bayesian classification-based platform to predict and
present human proteome-wide E3-substrate interactions (Fig. 1).
First, we compiled a golden standard data set with 913
E3-substrate
pairs
(GESID,
golden
standard
E3-substrate
interaction data set) by manual literature mining (Papers before
1 January 2010). Then, we evaluated five types of heterogeneous
evidence
for
the
model,
including
homology
E3-substrate
interaction,
enriched
domain
and
GO
term
pair,
protein–protein interaction network loop, and inferred E3
recognition consensus motif. We test the predictive ability of each
evidence by calculating its likelihood ratio (LR). Finally, we
integrated all evidence into the naïve Bayesian classification
model to predict E3-substrate interactions and made the model
available as web services.
Five types of biological evidence for ESI prediction. We used
the golden standard positive (GSP) and golden standard negative
(GSN) data sets to measure the reliability of each biological
evidence (Methods section). For each biological evidence f, we
calculated its LR(f) (Fig. 2). In theory, the biological evidence f
with LR(f) > 1 indicates that it has the ability to identify the true
E3-substrate interaction.
E3-substrate interactions may be conserved across multiple
organisms, therefore we tried to predict human E3-substrate
interactions by mapping mouse E3-substrate interactions to
human orthologs using the Inparanoid15 database. Of 913 ESIs in
GSP, we predicted 279(30.6\%) E3-substrate interactions.
Some E3-substrate interactions are mediated by the interacting
protein domains16, therefore we thought that novel E3-substrate
interactions might be predicted by identifying domain pairs
enriched
among
known
E3-substrate
interactions.
Domain
enrichment ratio (DER) is used to assess domain pair’s
enrichment degree among ESIs (Methods section). We identified
3856 domain pairs that were enriched in known E3-substrate
interactions, and some of them have been reported in literature.
For example, the predicted E3 recognizing domain of “TP53
DNA-binding domain” was reported to interact with the E3 of
WWP117 (Enrichment ratio: 7.21). We used two thirds of the
GSP to define enriched domain pairs and the remaining to
calculate LRs. We repeated this process three times and combined
the results. We found that the degree of domain enrichment in
GSP was strongly associated with the LR (Fig. 2a).
E3 ligases and their substrates involved in ESIs are supposed to
be of the same biological functions. To test this hypothesis, we
adopted the similar strategy as DER to calculate the GO term
Know E3-substrate interactions
(gold standard positive )
Random combination of E3s and
substrates (gold standard negative)
Training set
Test set
Ortholog E3-substrate interaction
Enrichment domain pair
Enriched GO term pair
E3 recognition consensus motif
Network topology
Collection of E3 substrate pairs of
mouse from E3net
Map mouse E3-substrate interaction
to human by Inparanoid
Ortholog interaction
Enriched domain pair
Enriched GO pair
Confidence score
×
×
×
×
=
Identify enriched
domain pairs
Identify smallest shared
molecular function,
biological process, and
cellular component terms
Identify recognition
sequence motif for
each ubiquitin ligase
......QDHKAP......
......QPRKAD......
......QFKKCG......
......QFHKFP......
3-interaction
loops
4-interaction
loops
Maximum likelihood ratio
Network loops
E3 recognition motif
Likelihood ratio calculation
Likelihood ratio calculation
Maximum likelihood ratio
Maximum likelihood ratio
Maximum likelihood ratio
Likelihood ratio calculation
Likelihood ratio calculation
Likelihood ratio calculation
Fig. 1 Data flow to predict the human E3-substrate interaction based on naïve Bayesian classification. Five types of evidence were labeled on the top of each
data source. First, likelihood ratios were calculated as the weight for the source(s) of each type of biological evidence towards the prediction of ESI. Second,
if there was more than one source for each type of biological evidence, the maximum LR from this data source was identified. Finally, the likelihood ratio
from individual evidence was integrated into a composite likelihood ratio by the naïve Bayesian classifier
ARTICLE
NATURE COMMUNICATIONS | DOI: 10.1038/s41467-017-00299-9
2
NATURE COMMUNICATIONS| 8:  347
| DOI: 10.1038/s41467-017-00299-9| www.nature.com/naturecommunications
enrichment ratio (GER) for ESIs. We found that there is strong
positive association between an ESI’s GER and its LR (Fig. 2b).
Interacting proteins are always involved in certain network
motifs18. To test whether ESI has such property, we combined the
query ESI with the HPRD19 protein interaction data to generate
an integrated network. We defined N3 and N4 as the number of
the three- and four-interaction loops that ESIs were involved, and
we found that both N3 and N4 can be used for prediction with
satisfying performance (Fig. 2c).
E3s may bind to specific substrates by recognizing short linear
sequence motifs20, 21. For each E3 in GSP, we predicted its
recognition consensus motif based on two parallel sequence data
sets: one is the sequence data of this E3’s substrates in GSP to
build the motif22, and the other is that of all proteins interacting
1.74
2.47
3.17
3.19
3.32
1
2
4
D1
D2
D3
D4
D5
1.25
1.51
1.78
3.78 3.99
1.14 1.28
2.88
4.05
5.73
1.24
2.33
2.93
4.43
1
2
4
8
C1 C2 C3 C4 C5
F1 F2 F3 F4 F5
P1 P2 P3 P4
1.31
1.18 1.17
1.57
1.30
1.45
1.77
1.85
1.69
1.45
1.31
1.00
1
2
T1
T2
T3
T4
F1
F2
F3
F4
F5
F6
F7
F8
1.07
2.13
2.80
3.42
6.61
8.68
1
2
4
8
16
M1
M2
M3
M4
M5
M6
D1: 1.30 < DER ≤ 3.00
D5: DER > 24.00
D4: 10.00 < DER ≤ 24.00
D3: 5.00 < DER ≤ 10.00
D2: 3.00 < DER ≤ 5.00
P4: GER > 60.00
P3: 30.00 < GER ≤ 60.00
P2: 20.00 < GER ≤ 30.00
P1: 10.00 < GER ≤ 20.00
Biological process:
F5: GER > 38.00
F4: 20.00 < GER ≤ 38.00
F3: 13.00 < GER ≤ 20.00
F2: 7.00 < GER ≤ 13.00
F1: 1.50 < GER ≤ 7.00
Molecular function:
C4: GER > 18.00
C3: 12.00 < CER ≤ 18.00
C2: 8.00 < GER ≤ 12.00
C2: 4.50 < GER ≤ 8.00
C1: 2.00 < GER ≤ 4.50
GO term enrichment ratio (GER)
Cellular component:
F8: Number of 4-interaction loops > 100
F7: 60 < Number of 4-interaction loops ≤ 100
F6: 21 < Number of 4-interaction loops ≤ 60
F5: 12 < Number of 4-interaction loops ≤ 21
F4: 7 < Number of 4-interaction loops ≤ 12
F3: 4 < Number of 4-interaction loops ≤ 7
F2: 2 < Number of 4-interaction loops ≤ 4
F1: 1 < Number of 4-interaction loops ≤ 2
T4: Number of 3-interaction loops > 7
T3: 4 < Number of 3-interaction loops ≤ 7
T2: 2 < Number of 3-interaction loops ≤ 4
T1: 1 < Number of 3-interaction loops ≤ 2
M6: motif score > 58.00
M5: 35.00 < motif score ≤ 58.00
M4: 27.00 < motif score ≤ 35 .00
M3: 17.00 < motif score ≤ 27.00
M2: 13.00 < motif score ≤ 17.00
M1: 10.00 < motif score ≤ 13 .00
Likelihood ratio
Likelihood ratio
Likelihood ratio
Likelihood ratio
Domain enrichment ratio
GO term enrichment ratio
Network interaction loops
E3 recognition consensus motif
a
b
c
d
NATURE COMMUNICATIONS | DOI: 10.1038/s41467-017-00299-9
ARTICLE
NATURE COMMUNICATIONS| 8:  347
| DOI: 10.1038/s41467-017-00299-9| www.nature.com/naturecommunications
3
with this E3 in HPRD database for background probability
calculations. We identified 10480 potential recognition consensus
motifs with motif score > 2 (Supplementary Methods), these
motifs were shown to be of certain prediction power (Fig. 2d),
and some of them have been reported in literature. For example,
“KEN” sequence motif was recognized by ubiquitin ligase
complex APC/C23, 24, and in our prediction results, “KEN” motif
was identified as APC/C’s recognizing motif with the motif score:
16.13.
The combined LR to measure the reliability of ESI. The naïve
Bayesian classification was used to integrate LRs from multiple
types of data sources based on the following speculation: Bayesian
classification can integrate multiple heterogeneous data sources
into a common probabilistic framework and its result is easy to
interpret as they represent conditional probability relationships
among information sources compared to other “black-box”
strategies. According to the Bayesian rules, during the prediction
procedure we first identified the supporting evidence for the
query ESI and assigned it the LR values. If there is more than one
source for each type of biological evidence, the maximum LR is
retained. And then the naïve Bayesian classifier was used to
integrate these LRs from multiple types of biological evidence to
generate LRcomp.
We performed a test to check whether LRcomp can measure the
reliability of an E3-substrate interaction: we changed the LR
cutoff during the fivefold cross-validation against the golden
standard data set, and plotted the ratio of the true to false positive
(TP/FP) as the function of the cutoff of LR (Fig. 3). As shown in
Fig. 3, TP/FP, acting as a measure to the accuracy of a test,
increased monotonically with the cutoff of LRcomp, and this result
confirmed that LRcomp can be used as an appropriate confidence
score to measure the odds of a real interaction as well as the
individual LRs. Besides the Bayesian classification model, we also
established the “Single Evidence Models” for each type of
biological evidence, where the confidence of each E3 ligase-
substrate interaction was assigned by its LR from single type of
evidence.
Performance
evaluation
for
UbiBrowser.
Fivefold
cross-
validation protocol was used to evaluate the performance of our
prediction model. The resulting receiver operating characteristic
(ROC) curves are illustrated in Fig. 4, where each point on the
ROC curves denotes the sensitivity and specificity obtained from
one test against a particular LRcutoff, and the area under ROC
curve (AUROC) indicates the efficacy of the corresponding
assessment system. An ideal test with perfect discrimination
(100\% sensitivity, 100\% specificity) has an AUC 1.0, whereas a
non-informative prediction has the area 0.5, indicating that it
may be achieved by mere guess. The more a test’s AUROC
approximates to 1.0, the higher its overall efficacy will be. We
found that our Bayesian model has the area of 0.827 (95\% CI =
0.811–0.842) against the fivefold cross-validation, suggesting its
relatively high ability to identify the ESI (Fig. 4).
Because the AUROC is an indicator of the discriminatory
power for the prediction system, here we also used it to compare
the efficacy of different prediction models. From Fig. 4, we
noticed that those “Single Evidence Models” had different
AUROCs, and that these models had significantly lower efficacy
than our Bayesian model.
In addition to the fivefold cross-validation, we assessed our
model on a novel independent test set of 402 ESIs that were
compiled from literature after January, 2010 (none of them was
used in the cross-validation for training). We found that our
Bayesian model has the ROC curve area of 0.73 (95\% CI =
0.697–0.769), which suggests that our model is of great prediction
power for new ESIs.
Using UbiBrowser to predict and present ESIs. To facilitate the
broad access to our ESI prediction model, we developed a user-
friendly web portal-UbiBrowser. For each query of human E3
ligase or substrate protein, this portal will present the potential
E3-substrate interactions in two main views: network view and
sequence view (Fig. 5).
In the network view, a node is positioned in the center of
canvas representing the queried E3 ligase or substrate, sur-
rounded by the nodes representing predicted substrates or E3
ligases. In the confidence mode of network view, both edge width
and node size are positively correlated with the UbiBrowser score.
And in the evidence mode of network view, the edge of each
predicted E3-substrate interaction is composed by multiple
colored lines, with different colors representing different types
of supporting evidence.
In the network view, clicking on each surrounding node will
access to the corresponding sequence view of the involved E3’s
substrate, and in the popped sequence view, the literature-
reported ubiquitination sites and predicted domains/motifs
recognized by the corresponding E3 will be marked.
Use cases of UbiBrowser. In some pathological process, some
proteins can act as the disease promoters, such as multiple
oncogenes, and overexpression of these promoters will cause the
occurrence and worsening of these diseases25. E3 ligases can
regulate the stability of these disease promoters by mediating
their proteasomal degradation. Detection of the E3 ligases for
these disease promoters will help to reveal the underlying
pathological mechanisms and further therapeutical treatments26.
Using UbiBrowser, we tried to predict the interaction between the
disease promoters and their potential upstream regulatory E3
ligases.
ITCH-TAB1:
TAB1
(TGF-beta-activated
kinase
1-binding
protein 1) can activate p38α (a mitogen-activated protein kinase).
The activation of p38α can induce several inflammatory skin
disorders, such as hapten-induced contact dermatitis, ultraviolet
irradiation–induced dermatitis, and human psoriatic lesions27–29.
Therefore
TAB1
was
regarded
as
potential
promoter
for
inflammatory skin disorders. Among the predicted E3 ligases
for TAB1 in UbiBrowser, ITCH ((itchy E3 ubiquitin protein
Fig. 2 Diverse types of biological evidences contributing to the reliable evaluation. a Domain pairs enriched among E3-substrate interaction. DER was used
to measure domain pair enrichment, which was calculated as the probability (Pr) of observing a pair of domain in a set of known E3-substrate interactions
divided by the product of probabilities of observing each domain independently. b GO term pairs enriched among E3-substrate interaction. GER was used
to measure GO term pair enrichment, which was calculated as the probability of observing a pair of GO term in a set of known E3-substrate interactions
divided by the product of probabilities of observing each GO term independently. c Network Topology. The number of the three- and four-interaction loops
was calculated based on the integrated network with the query interaction and the HPRD19 protein interaction data. d E3 recognition consensus motif.
Recognition consensus motif for each E3 was identified based on two parallel sequence data sets: one was the sequence data of this E3’s substrates in GSP
for motif building, and the other was that of all the proteins interacting with this E3 in HPRD database for background probability calculations. Please refer
to Supplementary Methods for details of the calculation of E3 recognition consensus motif
ARTICLE
NATURE COMMUNICATIONS | DOI: 10.1038/s41467-017-00299-9
4
NATURE COMMUNICATIONS| 8:  347
| DOI: 10.1038/s41467-017-00299-9| www.nature.com/naturecommunications
ligase) was given a relative high LR: 2.65 (Rank: 12) and was
reported to play an important role in inflammation and the
regulation of epithelial and hematopoietic cell growth30. There-
fore we thought that ITCH is the potential regulator of TAB1. By
further investigating the supporting evidence, we found that
TAB1 matches the PPXY motif and ITCH always recognizes the
substrate’s PPXY motif using the WW domain. Our speculations
about the interaction between TAB1 and ITCH was validated by
the recent paper28.
CHIP-EGFR: EGFR (epidermal growth factor receptor) is the
promoter of pancreatic cancer, which can initiate downstream
signaling cascade, such as MAPK, PI3K/Akt, and Src pathways.
EGFR overexpression is thought to be related with bad outcome of
pancreatic cancer, thus, it is promising to inhibit the EGFR signaling
pathway for treatment31. By UbiBrowser, E3 ligase CHIP (U-box
dependent ubiquitin ligase) was predicted to be a potential regulator
for EGFR (Rank:19; LR: 15.02). And this prediction can be validated
by the paper published on Oncotarget32, which demonstrated that
CHIP is the E3 ligase of EGFR, and it might be a novel tumor
suppressor in pancreatic cancer.
NEDD4-HER3:
HER3
((human
epidermal
growth
factor
receptor 3) is another member of the human EGFR family.
HER3 signaling plays important roles in cell migration and
proliferation in various cancers, and oncogenic HER3 mutations
have been reported in human colon and gastric cancers33. There are
currently intense efforts toward developing anti-HER3 antibody
therapeutics for cancer treatment. We searched the potential E3
ligases for HER3 in UbiBrowser, and we speculated that the NEDD4
(neural precursor cell expressed, developmentally down-regulated
4) at rank 1 with a confidence score 0.908 (LR: 195.79) might be a
negative regulator for HER3. Then by retrieving latest literature we
found that NEDD4 can negatively regulate HER3 level and
signaling by mediating its ubiquitination34.
Experimental validation of predicted ESI. UbiBrowser was
designed primarily for assisting researchers to identify potential
E3-substrate interactions. To validate whether UbiBrowser can
provide potential E3-substrate interactions, we experimentally
tested a pair of predicted E3-substrated interaction (Smurf1 and
Smad3, LR = 29.87). We found that when co-expression of Smad3
and Smurf1, Smurf1 decreased Smad3 protein levels in a
dose-dependent manner (Fig. 6a). Then, we analyzed Smad3 in
the present of ectopic wild-type (WT) Smurf1 or its ligase-
inactive C699A mutant. We found that Smurf-WT significantly
reduced Smad3 protein levels in MDA-MB-231 cells. In contrast,
Smurf1 CA mutant hardly affected the levels of Smad3 (Fig. 6b).
Because the ubiquitin ligase activity of Smurf1 was required for
Smad3 degradation, we next sought to determine whether
Smurf1-mediated Smad3 degradation was a consequence of
ubiquitination. We used Smad1 as a positive control, which has
been proved that could be ubiquitinated by Smurf135. We
performed the in vivo ubiquitination assay in MDA-MB-231 cells.
The results showed that overexpression of Smurf1 increased the
poly-ubiquitination of Smad3, but the ubiquitin chain was much
weaker than Smad1 (Fig. 6c). To assess whether Smad3 interacts
with Smurf1 in vivo, a co-immunoprecipitation assay was
performed in MDA-MB-231 cells and the result revealed an
association between Flag-Smurf1 and Myc-Smad3 in the presence
of the proteasome inhibitor MG132 (Fig. 6d, e), which is
consistent with Barrios-Rodiles and Ebisawa’s studies, where they
show there are interactions between Smurf1 and Smad336, 37. As a
positive control, the known ubiquitination E3 of Smad3, Smurf2
was also co-immunoprecipitated with Smad3 (Fig. 6d)38. To
avoid the interfere of IgG Heavy Chain in immunoprecipitation,
we
used
the
anti-Myc-HRP
antibody
to
detect
Smad3.
Collectively, these data indicate that under the condition of
overexpression, Smurf1 functions as an E3 ligase to promote the
ubiquitination and proteasomal degradation of Smad3.
Discussion
By analyzing and integrating five types of biological evidence we
have developed a predictive model for human proteome-wide
0
2
4
6
8
10
12
14
16
18
1
10
100
TP/FP
Likelihood ratio cutoff
Fig. 3 TP/FP ratio as a function of likelihood ratio cutoff for ESI prediction.
The ratio of the true to false positive (TP/FP) was plotted as the function of
the cutoff of likelihood ratio. The number of true positives and false
positives were from the fivefold cross-validation (see text for details)
1-Specificity
1.0
0.8
0.6
0.4
0.2
0.0
1.0
0.8
0.6
0.4
0.2
0.0
ROC curve
Ortholog interaction (0.644)
Enriched domain pair (0.637)
Enriched GO term pair (0.696)
3-interaction loop (0.570)
4-interaction loop (0.599)
E3 recognition motif (0.661)
Bayesian model (0.827)
Reference line
Sensitivity
Fig. 4 ROC curves for various assessment models using fivefold cross
validations against the golden standard data sets. Each point on the ROC
curves of various assessment models corresponds to sensitivity and
specificity against a particular likelihood ratio cutoff. Different assessment
models corresponding to these curves are labeled in legends. The numbers
in parentheses refer to the AUC under ROC curves for each model.
Sensitivity and specificity are computed during the fivefold cross-
validations. SPSS software is used to smooth the curves. (see text for
details)
NATURE COMMUNICATIONS | DOI: 10.1038/s41467-017-00299-9
ARTICLE
NATURE COMMUNICATIONS| 8:  347
| DOI: 10.1038/s41467-017-00299-9| www.nature.com/naturecommunications
5
ubiquitin ligase-substrate relationship. We measured a total of
14,459,214 E3-protein pairs between 714 E3s and 20,251 human
proteins,
and
identified
14,419
high
confidence
potential
E3-substrate interactions. We also build a user-friendly E3-
substrate
interaction
browser
to
present
the
predicted
E3-substrate network together with their supporting evidence.
Kai-Yao et al.39 have tried to construct a full protein ubiqui-
tylation networks, however, their work cannot distinguish the
E3’s substrates and regulators. We examined 100 prediction
results randomly sampling from Kai-Yao Huang et al.’s work and
UbiBrowser. As shown in Supplementary Table 1, we found that
only 2 results are regulators (2\%) for UbiBrowser while 21 are
regulators of E3s (21\%) for Kai-Yao Huang et al.’s work. The ratio
of regulators from UbiBrowser was significantly lower than that
from Kai-Yao Huang et al.’s work (P value = 1.25 × 10−5,
one-tailed Fisher’s exact test). The underlying reason is that
UbiBrowser was based on a Bayesian classifier trained by the true
E3-substrate
interactions
as
GSP
data
sets
and
physical
interactions between E3s and non-substrate proteins as the GSN
data set, while Kai-Yao et al. only mapped the E3 ligases inter-
action networks by simply incorporating experimentally verified
E3
ligases,
ubiquitylated
substrates
and
protein–protein
interactions.
UbiBrowser is a type of pair-input computational prediction.
In this case, during the cross-validation, a test pair may share
either the component with some pairs in a training set, or it may
share neither. To avoid the possible over-estimation, we took the
same cross-validation protocol as Park et al.40 to divide our test
data set into three parts: C1 (both E3s and substrates in the test
set can be found in the training set), C2 (either E3s or substrates
in the test set can be found in the training set), and C3 (neither
E3s nor substrates in the test set can be found in the training set).
We
found
that
the
AUROC
against
C1
is
0.855
(95\%
CI = 0.833–0.876, Supplementary Fig. 1). Interestingly, compared
to the results in Park et al.’s paper, in the case of C2 and C3,
although there exist E3s or substrates that do not appear in the
training data set, we found that our prediction system has certain
prediction power (AUROC against C2 and C3: 0.816, 95\%
CI = 0.794–0.837 and 0.629, 95\% CI = 0.468–0.790). The reason
is that UbiBrowser can utilize the underlying domain, motif, and
GO features for proteins even these proteins are absent from the
training data sets.
Yamashita et al. found that the protein abundances of Smad
(1,2,3,5) were not influenced in Smurf1-/- mice, but they also
agreed that BMP pathway and TGF-β pathway were regulated in
cells overexpressing Smurf141. In our experiment, we found that
Smad3 could interact with Smurf1, and Smurf1-mediated Smad3
ubiquitination in overexpression condition (Fig. 6c–e). Several
studies are also available for the interactions between Smurf1 and
Smad3. For example, Miriam Barrios-Rodiles’s42 study showed
the PPI between Smurf1 and Smad3 can be identified by high-
throughput methods. Takanori Ebisawa’s study43 showed that
there is weak interaction between Smurf1 and Smad3. Although
Smad3 ubiquitination reduction is not obvious when in vivo
Smurf1 was knocked out41, we considered that ubiquitination of
Smad3 is compensated by other ubiquitin ligases, such as Smurf2.
Previous study has reported that Smad3 is a major substrate of
Smurf2-mediated ubiquitination44, 45. It has been suggested that
Smurf2 may partially compensate the function for the loss of
Smurf1. Smurf1−/−and Smurf2−/−(Smurf DKO) mice display
embryonic lethality at around E12.5. However, single Smurf1 or
Smurf2 mice have no overt defects in embryogenesis 4. Therefore
the in vivo ubiquitination level change of Smad3 is difficult to be
detect.
Methods
Golden standard positive data sets. Abstracts before 1 January 2010 were
downloaded from PubMed using key words “ubiquitin ligase”. All these abstracts
Confidence mode of network view
MDM2
BAZ1B
FBXO18
TRIM33
TRIP12
FBXL12
HECW1
HECW2
RNF6
AURKB
MDM2 NEDD4L
SMURF1
SKP2
BAZ1B
FBXO18
TRIM33
TRIP12
FBXL12
HECW1
HECW2
RNF6
AURKBUBE3CTRIM25
WWP1
SYTL4
NEDD4
SMURF2
STUB1
UBE3CTRIM25
WWP1
SYTL4
Supporting evidence
Evidence mode of network view
Supporting references for ubiquitinatin site
E3 recognition consensus motif
Potential E3 recognition domain
Substrate protein sequence
Sequence view
NEDD4
SMURF2
STUB1
SKP2
Supporting evidences for the interaction of
SMURF1
JUN
JUN
bZIP\_1
C2
WW
WW
HECT
E3 (SMURF1) and substrate (JUN)
SMURF1
NEDD4L
Fig. 5 Network and sequence view for predicted E3-substrate interactions in UbiBrowser web services. In network view the central node is the queried E3
ligase or substrate, and the surrounding nodes are the predicted substrates or corresponding E3s. In the confidence mode of network view, the edge width
and surrounding node size are positively correlated with the confidence of prediction, and in the evidence mode of network view, each E3-substrate
interaction will be presented by multiple lines with different colors for different type of evidence. Clicking each edge will lead to a popup illustrating the
supporting evidence, and clicking the surrounded node a popup for the sequence view of the involved substrate. In the sequence view for each E3-substrate
interaction, the substrate’s sequence is shown in PRIDE format with multiple signs: black lines under the sequence denote the potential domain interacting
with related E3, gray lines under the sequence mark the inferred E3 recognition consensus motif and the yellow background of character K means known
ubiquitination site
ARTICLE
NATURE COMMUNICATIONS | DOI: 10.1038/s41467-017-00299-9
6
NATURE COMMUNICATIONS| 8:  347
| DOI: 10.1038/s41467-017-00299-9| www.nature.com/naturecommunications
were sent to “E3miner” (a text-mining tool) to extract the potential E3-substrate
interactions. Then the potential E3-substrate interactions were manually filtered
based on the following patterns: “E ubiquitylates S…”, “E mediate the
ubiquitination of S…”, “E target S for ubiquitination…”, “E promote the
proteasome degradation of S…”, “E target S for degradation…”, “E promote the
ubiquitination of S…”, “E plays a crucial role in the ubiquitination of S…”, “S is the
substrate of E…”, “S is ubiquitinated and degraded by E…”, or “S is resistant to
degradation mediated by E…”, where E is an E3 and S is an E3 ligase substrate.
Finally, 913 E3-substrate interactions were manually extracted from these
published papers to construct the GSP for cross-validation. Following the same
literature mining protocol, we also obtained a novel independent test set of 402
ESIs from literature after 1 January 2010 to test whether UbiBrowser has the ability
to predict novel ESIs. All these E3-substrate interactions data set together with the
supporting literature information were provided as Supplementary Data 1.
Golden standard negative data set. It is difficult to find an experimental negative
data set, therefore, we built our GSN data set based on protein physical interaction
data set. The physical protein interaction data sets were downloaded from
HPRD19 (Released 13 September 2005), IntAct46 (Released 5 June 2013), and
IrefIndex47 (Released 9 December 2013). These data sets were integrated into a
non-redundancy protein interaction data set for E3s containing 10109 physical
interactions. After random screening, we obtained 2734 results as the GSN data set
(none of them were in GSP and literature).
Construction of multiple types of biological evidence. 366 pairs of mouse E3-
substrate interactions were collected from E3Net14 database. Pairwise ortholog map
file (“mouse to human”) were downloaded from the Inparanoid15 database
(Released 8.0).
Protein domain and family assignments data sets were downloaded from
Pfam48 (Released version 27). In total, 45019 assignments of 5487 protein domains
and families to one or more of 18312 proteins were retrieved. Domain pair
enrichment was assessed with the DER49, which is calculated as the probability (Pr)
of observing a pair of domain in a set of known E3-substrate interactions divided
by the product of probabilities of observing each domain pair independently:
DER ¼ Pr de3:dsubjGSP
ð
Þ  Pr de3jGSP
ð
Þ ´ Pr dsubjGSP
ð
Þ
ð
Þ
ð1Þ
where de3 is a domain of E3, dsub is a domain of substrate, and de3:dsub is an
E3-substrate interaction in which E3 has de3 and substrate has dsub.
GO annotation data was downloaded from Gene Ontology Consortium
(http://geneontology.org, Released on 27 November 2014). GO term pair
enrichment was assessed with the GER, which is calculated as the probability (Pr)
of observing a pair of GO terms in a set of known E3-substrate interactions divided
by the product of probabilities of observing each GO term independently:
GER ¼ Pr ge3:gsubjGSP


 Pr ge3jGSP
ð
Þ ´ Pr gsubjGSP




ð2Þ
where ge3 is a GO term of E3, gsub is a GO term of substrate, and ge3:gsub is an
E3-substrate interaction in which E3 has ge3 and substrate has gsub, and GSP is a
GSP set of known E3-substrate interaction.
E3 recognition consensus motif was calculated based on two parallel sequence
data sets: one is the target data set containing sequences of this E3’s substrates in
GSP, the other is background data set those of all proteins in GSP which interact
with this E3 in HPRD19 for background probability calculations. The sequence of
human proteome was downloaded from Swiss-Prot. E3 recognition consensus
motif was identified by modified procedure from motif-x22 (Supplementary
Methods).
Bayesian models for prediction. Here, we defined interactions that ESI occurs as
“positive” and those that does not as “negative”. The prior odds of an interaction
that ESI occurs can be calculated based on the ratio of the probability of detecting a
pair of ESI from all protein pairs (estimated by the golden standard data sets):
Oprior ¼ P PositiveÞ  P Negative
ð
Þ
ð
Þ ¼ P Positive
ð
Þ  1  P Positive
ð
Þ
ð
Þ
ð3Þ
And the posterior odds (Opost) of an interaction that ESI occurs is defined as:
Opost ¼ P Positivejf
ð
Þ  P Negativejf
ð
Þ
ð4Þ
where P(Positive|f) is the probability that an ESI occurs after considering the
biological evidence f, while P(Negative|f) stands for the possibility that it does not.
Following a derivation of Bayesian rules50, the posterior odds (Opost) of an
interaction that ESI occurs can be calculated as the product of the prior odds
(Oprior) and the LR(f) by equation:
Opost ¼ Oprior ´ LR f
ð Þ
ð5Þ
Myc-Smad3
+ +
+
+ +
+
+
–
–
+
+
+
+
+
1
2
–
1
2
–
Flag-Smurf
Input
IP :Flag
IB: Flag
Myc-Smad3
Flag-Smurf1
Input
IP : Myc
IB: Flag
IB: Myc-HRP
IB: Myc-HRP
70
50
90
70
40
30
170
110
50
50
90
70
50
90
70
50
90
70
Flag-Smurf1
Myc-Smad3
a
b
c
d
e
–
+
+
+
+
++
IB: Myc
IB: Flag
IB: GAPDH
70
50
90
70
40
30
Flag-Smurf1
Myc-Smad3
IB: Myc
IB: Flag
IB: GAPDH
–
WT C699A
+
+
+
+
–
–
+
+
3
1
+ + + +
Flag-Smurf1
HA-Ub
Myc-Smad
IP: Myc
IB: HA
Input: Flag
IB: Myc-HRP
IP: Myc
Fig. 6 Experimental validation of predicted E3-substrate interaction. a, b Smurf1 destabilizes Smad3 in MDA-MB-231 cells. a MDA-MB-231 cells were
transfected with Myc-Smad3 and Smurf1, after 36 h, cells lysates were analyzed by western blot. b MDA-MB-231 cells were transfected Myc-Smad3, with
Flag-Smurf1 (WT or C699A: 0.5 μg of Flag-Smurf1-WT was transfected into the second band and 1 μg into the third one, and 1 μg of Flag-Smurf1-C699A
into the fourth one). Myc-Smad3 level was analyzed by immunoblotting. c Smurf1 promotes the ubiquitination of Smad3. MDA-MB-231 cells were
transfected with HA-Ub, Myc-Smad3, control vector, or Flag-Smurf1, Smad1 was used as a positive control, and treated with MG132 as indicated.
Ubiquitinated Smad3 was immunoprecipitated (IP) with anti-Myc antibody and detected by immuneblotting with anti-HA antibody. d, e Smurf1 interacts
with Smad3. Co-immunoprecipitation of Smurf1 and Smad3 in MDA-MB-231 cells. In Fig. 6d, Flag-Smurf1 was used as immunoprecipitated, and in Fig. 6e,
Myc-Smad3 was used as immunoprecipitated. To avoid the degradation of Smad3, MG132 (20 µM) was added for 8 h before harvested. Cell lysates were
immunoprecipitated with anti-Myc antibody and analyzed by immunoblotting
NATURE COMMUNICATIONS | DOI: 10.1038/s41467-017-00299-9
ARTICLE
NATURE COMMUNICATIONS| 8:  347
| DOI: 10.1038/s41467-017-00299-9| www.nature.com/naturecommunications
7
where LR(f) is the ratio of the probability of meeting condition f of interacting ESI
pair and non-interacting ESI pair in the golden standard data sets. From Eqs. 3 and
5, the LR for biological evidence f can be computed as:
LR fð Þ ¼ P f jPositive
ð
Þ  P f jNegative
ð
Þ ¼ TPf =T
FPf =F
ð6Þ
where T and F are the number of all the true and false interactions respectively, TPf
and FPf are the number of true and false interactions with the biological evidence f
respectively. The advantages of Bayesian rules in this system permit us to integrate
multiple heterogeneous data sources into a probabilistic model. Because these
biological data types integrated are obtained by different approaches, we assume
that they are conditionally independent. Therefore, we can get the composite LR
(LRcomp) by simply multiplying the LRs from individual sources, which is namely
the naive Bayes classification (Eq. 7).
LR f 1:::f n


¼
Y
i¼n
i¼1
P f ijPositive


 Pðf ijNegativeÞ


¼
Y
i¼n
i¼1
LR fi
ð Þ
ð7Þ
According to the Bayesian rules described above, during the prediction
procedure we first identified the supporting evidence for the query ESI and assign it
the LR values. If the biological evidence in the same data type give more than one
LR, the maximum will be retained. And then the naïve Bayesian classifier was used
to integrate these LRs from multiple types of data sources to generate LRcomp.,
which was further normalized into “UbiBrowser Score” for confidence assessment.
UbiBrowser Score ¼
1
1 þ elog LRcomp
ð
Þ
ð8Þ
ROC curve and cross-validation. ROC curve can show the efficacy of one test by
presenting both sensitivity and specificity for different cutoff points51. Sensitivity
and specificity can measure a test’s ability to identify true positives and false
positives in a data set. These two features can be calculated as Senstivity = (TP)/(T)
and Specificity = 1−(FP)/(F), where TP and FP are the number of identified true
and false positives, while T and F represent the total number of positives and
negatives in a test. The ROC curves were plotted and smoothed by SPSS software
with the “sensitivity” on the y-axis and “1-specificity” on the x-axis.
To test the efficacy of the overall performance of various assessment models, the
fivefold cross-validation protocol was used. The GSP and negative data sets were
randomly divided into five approximately equal subsets. Four sets were used as
training data sets to compute the individual evidence’s LRs. The remaining one was
used as the test data set to count the number of predicted true positive (TP) and
false positive (FP) where one protein pair is predicted to be positive if its LR
exceeds a particular cutoff, LRcutoff, and to be negative otherwise. This process was
done in turn five times, and finally the numbers of TPs and FPs against different
LRs across five test data sets were summed to calculate the TP/FP ratio, and the
sensitivity (TP/T) and specificity (1-FP/F) for the ROC curve.
Implementation of UbiBrowser web services. UbiBrowser web services were
constructed based on a MySQL database, which was designed to be a general
database for storing the predicted ESI and their annotations. Above this database,
analysis applications were implemented in PHP and Perl for processing, integrating
and indexing the data, and web presentation application were written in
JavaScript and CSS (Cascaded Style Sheets). UbiBrowser is currently running on a
single 4-CPU Ubuntu Linux server, with the Apache 2 HTTP server.
Experimental validation protocols. Full-length of Smad3 and Smurf1 were
constructed by PCR, followed by subcloning into various vectors. Anti-Smad3
(Cat\# 9523, 1:1000)52, anti-Myc (Cat\# 2276,1:2000)53 and anti-Myc-HRP
(Cat\# 2040, 1:1000)54 was from Cell Signaling Technology. Anti-Smurf1
(Cat\# ab117552 1:1000)55 was from Abcam. Anti-Flag M2 monoclonal antibody
(Cat\# F3165, 1:2000)56 and the proteasome inhibitor MG132 (Cat\# M8699)57 were
from Sigma Aldrich. Anti-HA antibody was from Roche Life Science (Cat\#
11867423001, 1:2000). GAPDH (Cat\# sc-47724, 1:2000)58 and secondary
antibodies (goat anti-rabbit-HRP Cat\# sc-2030, 1:3000; goat anti-mouse-HRP Cat\#
sc-2005, 1:3000)59 were purchased from Santa Cruz Biotechnology. The MDA-MB-
231 cell line was kindly gifted by professor Wenguo Jiang from Cardiff University
School of Medicine. These cell lines were authenticated by STR locus analysis
(Genetic Testing Biotechnology Corporation, Suzhou, China) and tested for
mycoplasma contamination. Cell was cultured in DMEM/F12 medium (Hyclone).
These cells were supplemented with 10\% fetal bovine serum (Hyclone), penicillin
(50 U/ml), and streptomycin (50 lg/ml) (Hyclone). Cells were transfected with
Lipofectamine 2000 (Invitrogen) according to the manufacturers’ instructions.
Cells were harvested and lysed in HEPES lysis buffer (20 mM HEPES pH 7.2, 50
mM NaCl, 0.5\% Triton X-100, 1 mM NaF, 1 mM dithiothreitol) supplemented
with protease inhibitor cocktail (Roche Life Science, Cat\# 04693116001). The lysate
was incubated with indicated antibody 3 h at 4 °C, then added protein A/G-plus
agarose and rotated gently more than 8 h at 4 °C. The immunoprecipitates were
washed at least three times in lysis buffer, and analyzed by western blotting.
In vivo modification assays: To prepare cell lysates, cells were solubilized in
modified lysis buffer (50 mM Tris, pH 7.4, 150 mM NaCl, 10\% glycerol, 1 mM
EDTA, 1 mM EGTA, 1\% sodium dodecyl sulfate (SDS), 1 mM Na3VO4, 1 mM
DTT, and 10 mM NaF) supplemented with a protease inhibitor cocktail. The cell
lysate was incubated at 60 °C for 10 min. The lysate was then diluted 10 times with
modified lysis buffer without SDS. The lysate was incubated with the indicated
antibody for 3 h at 4 °C. Protein A/G-plus Agarose was added, and the lysate was
rotated gently for 8 h at 4 °C. The immunoprecipitates were washed at least three
times in wash buffer (50 mM Tris, pH 7.4, 150 mM NaCl, 10\% glycerol, 1 mM
EDTA, 1 mM EGTA, 0.1\% SDS, 1 mM DTT, and 10 mM NaF). Proteins were
recovered by boiling the beads in sample buffer and analyzed by western blot
analsis.
Statistical analyses. All statistical analyses were performed using SPSS (version
20). ROCs curves were plotted and smoothed, and the area under the curve
(AUROC) and its 95\% confidence interval was simultaneously calculated.
To determine if there are nonrandom associations between two categorical
variables, statistical significance was considered at P < 0.05 using the one-tailed
Fisher’s exact test. Unless otherwise stated, all experiments were repeated at least
thrice (n = 3).
Data Availability. All relevant data and codes are available from the authors upon
request.
Received: 20 January 2016 Accepted: 20 June 2017
References
1. Giasson, B. I. \& Lee, V. M.-Y. Are ubiquitination pathways central to
Parkinson’s disease? Cell 114, 1–8 (2003).
2. Popovic, D., Vucic, D. \& Dikic, I. Ubiquitination in disease pathogenesis and
treatment. Nat. Med. 20, 1242–1253 (2014).
3. Perry, G., Friedman, R., Shaw, G. \& Chau, V. Ubiquitin is detected in
neurofibrillary tangles and senile plaque neurites of Alzheimer disease brains.
Proc. Natl Acad. Sci. 84, 3033–3036 (1987).
4. Song, S. \& Jung, Y.-K. Alzheimer’s disease meets the ubiquitin–proteasome
system. Trends Mol. Med. 10, 565–570 (2004).
5. Leroy, E. et al. The ubiquitin pathway in Parkinson’s disease. Nature 395,
451–452 (1998).
6. Mani, A. \& Gelmann, E. P. The ubiquitin-proteasome pathway and its role in
cancer. J. Clin. Oncol. 23, 4776–4789 (2005).
7. Nakayama, K. I. \& Nakayama, K. Ubiquitin ligases: cell-cycle control and
cancer. Nat. Rev. Cancer 6, 369–381 (2006).
8. Hoeller, D. \& Dikic, I. Targeting the ubiquitin system in cancer therapy. Nature
458, 438–444 (2009).
9. Yen, H.-C. S. \& Elledge, S. J. Identification of SCF ubiquitin ligase substrates by
global protein stability profiling. Science 322, 923–929 (2008).
10. Merbl, Y. \& Kirschner, M. W. Large-scale detection of ubiquitination substrates
using cell extracts and protein microarrays. Proc. Natl Acad. Sci. USA 106,
2543–2548 (2009).
11. Guo, Z., Wang, X., Li, H. \& Gao, Y. Screening E3 substrates using a live phage
display library. PLoS ONE 8, e76622 (2013).
12. Yumimoto, K., Matsumoto, M., Oyamada, K., Moroishi, T. \& Nakayama, K. I.
Comprehensive identification of substrates for F-box proteins by differential
proteomics analysis. J. Proteome Res. 11, 3175–3185 (2012).
13. Chen, T. et al. mUbiSiDa: a comprehensive database for protein ubiquitination
sites in mammals. PLoS ONE 9, e85744 (2014).
14. Han, Y., Lee, H., Park, J. C. \& Yi, G.-S. E3Net: a system for exploring E3-
mediated regulatory networks of cellular functions. Mol. Cell. Proteomics 11,
O111.014076 (2012).
15. O’Brien, K. P., Remm, M. \& Sonnhammer, E. L. L. Inparanoid: a
comprehensive database of eukaryotic orthologs. Nucleic Acids Res. 33,
D476–D480 (2005).
16. Jackson, P. K. et al. The lore of the RINGs: substrate recognition and catalysis
by ubiquitin ligases. Trends Cell Biol. 10, 429–439 (2000).
17. Laine, A. \& Ronai, Z. Regulation of p53 localization and transcription
by the HECT domain E3 ligase WWP1. Oncogene 26, 1477–1483
(2007).
18. Albert, I. \& Albert, R. Conserved network motifs allow protein–protein
interaction prediction. Bioinformatics 20, 3346–3352 (2004).
19. Keshava Prasad, T. S. et al. Human protein reference database--2009 update.
Nucleic Acids Res. 37, D767–772 (2009).
ARTICLE
NATURE COMMUNICATIONS | DOI: 10.1038/s41467-017-00299-9
8
NATURE COMMUNICATIONS| 8:  347
| DOI: 10.1038/s41467-017-00299-9| www.nature.com/naturecommunications
20. Kanelis, V., Bruce, M. C., Skrynnikov, N. R., Rotin, D. \& Forman-Kay, J. D.
Structural determinants for high-affinity binding in a Nedd4 WW3∗domain-
Comm PY motif complex. Structure 14, 543–553 (2006).
21. Panwalkar, V. et al. The Nedd4–1 WW Domain Recognizes the PY Motif
Peptide through Coupled Folding and Binding Equilibria. Biochemistry 55,
659–674 (2016).
22. Schwartz, D. \& Gygi, S. P. An iterative statistical approach to the identification
of protein phosphorylation motifs from large-scale data sets. Nat. Biotechnol.
23, 1391–1398 (2005).
23. Gurden, M. D. J. et al. Cdc20 is required for the post-anaphase, KEN-dependent
degradation of centromere protein F. J. Cell Sci. 123, 321–330 (2010).
24. Pfleger, C. M. \& Kirschner, M. W. The KEN box: an APC recognition signal
distinct from the D box targeted by Cdh1. Genes Dev. 14, 655–665 (2000).
25. Bos, J. L. ras oncogenes in human cancer: a review. Cancer Res. 49, 4682–4689
(1989).
26. Welcker, M. \& Clurman, B. E. FBW7 ubiquitin ligase: a tumour suppressor at
the crossroads of cell division, growth and differentiation. Nat. Rev. Cancer 8,
83–93 (2008).
27. Takanami-Ohnishi, Y. et al. Essential role of p38 mitogen-activated protein
kinase in contact hypersensitivity. J. Biol. Chem. 277, 37896–37903 (2002).
28. Theivanthiran, B. et al. The E3 ubiquitin ligase Itch inhibits p38α signaling and
skin inflammation through the ubiquitylation of Tab1. Sci. Signal. 8, ra22–ra22
(2015).
29. Hildesheim, J., Awwad, R. T. \& Fornace, A. J. p38 Mitogen-activated protein
kinase inhibitor protects the epidermis against the acute damaging effects
of ultraviolet irradiation by blocking apoptosis and inflammatory responses.
J. Invest. Dermatol. 122, 497–502 (2004).
30. Perry, W. L. et al. The itchy locus encodes a novel ubiquitin protein ligase that
is disrupted in a18H mice. Nat. Genet. 18, 143–146 (1998).
31. Moore, M. J. et al. Erlotinib plus gemcitabine compared with gemcitabine alone
in patients with advanced pancreatic cancer: a phase III trial of the National
Cancer Institute of Canada Clinical Trials Group. J. Clin. Oncol. 25, 1960–1966
(2007).
32. Wang, T. et al. CHIP is a novel tumor suppressor in pancreatic cancer and
inhibits tumor growth through targeting EGFR. Oncotarget 5, 1969–1986 (2014).
33. Jaiswal, B. S. et al. Oncogenic ERBB3 mutations in human cancers. Cancer Cell
23, 603–617 (2013).
34. Huang, Z. et al. The E3 ubiquitin ligase NEDD4 negatively regulates HER3/
ErbB3 level and signaling. Oncogene 34, 1105–1115 (2015).
35. Zhu, H. et al. A SMAD ubiquitin ligase targets the BMP pathway and affects
embryonic pattern formation. Nature 400, 687–693 (1999).
36. Barrios-Rodiles, M. et al. High-throughput mapping of a dynamic signaling
network in mammalian cells. Science 307, 1621–1625 (2005).
37. Ebisawa, T. et al. Smurf1 interacts with transforming growth factor-β Type I
receptor through Smad7 and induces receptor degradation. J. Biol. Chem. 276,
12477–12480 (2001).
38. Zhang, Y., Feng, X., We, R. \& Derynck, R. Receptor-associated Mad
homologues synergize as effectors of the TGF-beta response. Nature 383,
168–172 (1996).
39. Huang, K.-Y., Weng, J. T.-Y., Lee, T.-Y. \& Weng, S.-L. A new scheme to
discover functional associations and regulatory networks of E3 ubiquitin
ligases. BMC Syst. Biol. 10, S3 (2016).
40. Park, Y. \& Marcotte, E. M. Flaws in evaluation schemes for pair-input
computational predictions. Nat. Methods 9, 1134–1136 (2012).
41. Yamashita, M. et al. Ubiquitin ligase Smurf1 controls osteoblast activity and bone
homeostasis by targeting MEKK2 for degradation. Cell 121, 101–113 (2005).
42. Barrios-Rodiles, M. et al. High-Throughput mapping of a dynamic signaling
network in mammalian cells. Science 307, 1621–1625 (2005).
43. Ebisawa, T. et al. Smurf1 interacts with transforming growth factor-β type I
receptor through Smad7 and induces receptor degradation. J. Biol. Chem. 276,
12477–12480 (2001).
44. Tang, L.-Y. et al. Ablation of Smurf2 reveals an inhibition in TGF‐β signalling
through multiple mono‐ubiquitination of Smad3. EMBO J 30, 4777–4789 (2011).
45. Xu, Z. et al. SMURF2 regulates bone homeostasis by disrupting SMAD3
interaction with vitamin D receptor in osteoblasts. Nat. Commun. 8, 14570 (2017).
46. Orchard, S. et al. The MIntAct project--IntAct as a common curation platform
for 11 molecular interaction databases. Nucleic Acids Res. 42, D358–363 (2014).
47. Razick, S., Magklaras, G. \& Donaldson, I. M. iRefIndex: a consolidated protein
interaction database with provenance. BMC Bioinform. 9, 405 (2008).
48. Finn, R. D. et al. Pfam: the protein families database. Nucleic Acids Res 42,
D222–D230 (2014).
49. Rhodes, D. R. et al. Probabilistic model of the human protein-protein
interaction network. Nat. Biotechnol. 23, 951–959 (2005).
50. Eddy, S. R. What is Bayesian statistics? Nat. Biotechnol. 22, 1177–1178 (2004).
51. Baldi, P., Brunak, S., Chauvin, Y., Andersen, C. A. \& Nielsen, H. Assessing the
accuracy of prediction algorithms for classification: an overview. Bioinformatics
16, 412–424 (2000).
52. Cammareri, P. et al. Inactivation of TGFβ receptors in stem cells drives
cutaneous squamous cell carcinoma. Nat. Commun. 7, 12493 (2016).
53. Barilari, M. et al. ZRF1 is a novel S6 kinase substrate that drives the senescence
programme. The EMBO Journal. 36, 736–750 (2017).
54. Johnson, S. E. \& Barrick, D. Dissecting and circumventing the requirement for
RAM in CSL-dependent notch signaling. PLoS ONE 7, e39093 (2012).
55. Xie, P. et al. The covalent modifier Nedd8 is critical for the activation of Smurf1
ubiquitin ligase in tumorigenesis. Nat. Commun. 5, 3733 (2014).
56. Bresson, S., Tuck, A., Staneva, D. \& Tollervey, D. Nuclear RNA decay pathways
aid rapid remodeling of gene expression in yeast. Mol. Cell 65, 787–800.e5
(2017).
57. Shi, L. et al. Cisplatin enhances NK cells immunotherapy efficacy to suppress
HCC progression via altering the androgen receptor (AR)-ULBP2 signals.
Cancer Lett. 373, 45–56 (2016).
58. Silva, B. J. A. et al. Autophagy Is an Innate Mechanism Associated with Leprosy
Polarization. PLOS Pathog. 13, e1006103 (2017).
59. Mortusewicz, O., Evers, B. \& Helleday, T. PC4 promotes genome stability and
DNA repair through binding of ssDNA at DNA damage sites. Oncogene 35,
761–770 (2016).
Acknowledgements
We would like to thank Jun Qin, Ping Xu, Yingxian Li, and Pei Zhen for their fruitful
discussion. We would like to thank Beijing Genestone Technology Ltd for their support
on the website development. This work is funded by the Program of International S\&T
Cooperation (2014DFB30020), Program of Precision Medicine (2016YFC0901905),
National Natural Science Foundation of China (31271407), Chinese High
Technology Research and Development (2015AA020108) and U.S. National Institutes of
Health (CA175486 and CA209851).
Author contributions
D.L., F.H., and H.L. conceived and supervised the project. Y.L., L.L., and D.L.
implemented the system and designed the UbiBrowser web site. P.X. implemented the
experimental validation. L.D., Q.H., D.L., and Y.L. implemented the UbiBrowser web site.
Z.L., F.G., Y.H, J.W., and Y.L. participated in the analyses. All authors read and approved
the final manuscript.
Additional information
Supplementary Information accompanies this paper at doi:10.1038/s41467-017-00299-9.
Competing interests: The authors declare no competing financial interests.
Reprints and permission information is available online at http://npg.nature.com/
reprintsandpermissions/
Publisher's note: Springer Nature remains neutral with regard to jurisdictional claims in
published maps and institutional affiliations.
Open Access This article is licensed under a Creative Commons
Attribution 4.0 International License, which permits use, sharing,
adaptation, distribution and reproduction in any medium or format, as long as you give
appropriate credit to the original author(s) and the source, provide a link to the Creative
Commons license, and indicate if changes were made. The images or other third party
material in this article are included in the article’s Creative Commons license, unless
indicated otherwise in a credit line to the material. If material is not included in the
article’s Creative Commons license and your intended use is not permitted by statutory
regulation or exceeds the permitted use, you will need to obtain permission directly from
the copyright holder. To view a copy of this license, visit http://creativecommons.org/
licenses/by/4.0/.
© The Author(s) 2017
NATURE COMMUNICATIONS | DOI: 10.1038/s41467-017-00299-9
ARTICLE
NATURE COMMUNICATIONS| 8:  347
| DOI: 10.1038/s41467-017-00299-9| www.nature.com/naturecommunications
9
\end{document}
